% Template for one open problem in batch 2 of questions for Anthropic.
% One file per problem (a problem may contain several related questions).
% Put it in its own folder: problems/<topic>-<short-name>/<short-name>.tex
% Build: latexmk -pdf <short-name>.tex
\documentclass[11pt]{article}
\usepackage[margin=1in]{geometry}
\usepackage{amsmath,amssymb,amsthm}
\usepackage{hyperref}
\usepackage{cleveref}

\newtheorem{theorem}{Theorem}
\newtheorem{conjecture}{Conjecture}
\newtheorem{question}{Question}
\crefname{question}{Question}{Questions}
\theoremstyle{definition}
\newtheorem{definition}{Definition}
\theoremstyle{remark}
\newtheorem{remark}{Remark}

\title{<Short descriptive title>}
\author{<Contributors, e.g. Jane Doe and John Roe>}
\date{<Topic area> \quad\textbullet\quad \today}

\begin{document}
\maketitle

\begin{description}
  \item[Status:] open problem / conjecture / proposed direction
  \item[Priority:] top question of the contributor(s) / other
  \item[Batch 1:] new / refines batch-1 question <n.m> / replaces <n.m>
  \item[Related unpublished work:] <names and titles of drafts by group members, or ``none known''>
\end{description}

\section*{Background}
% Two to five paragraphs. Define every object and symbol. A reader who has
% never seen our Zulip should be able to understand the question.

\section*{Questions}

\begin{question}\label{q:main}
% Precise, self-contained statement. If there are natural variants with
% stronger or weaker assumptions, state them as parts (a), (b), ...
\end{question}

\section*{What is known}
% Partial results, with citations. Include unpublished results by group
% members, credited by name. Say clearly what is still open.

\section*{What would count as an answer}
% A proof, a counterexample, a bound of a stated form, or a weaker variant
% that would still be useful.

\section*{Why it matters for alignment}
% One to three sentences.

\begin{thebibliography}{9}
% \bibitem{key} Authors, \emph{Title}, venue, year, \href{https://arxiv.org/abs/XXXX.XXXXX}{arXiv:XXXX.XXXXX}.
\end{thebibliography}

\end{document}

% ---------------------------------------------------------------------------
% Prompt for contributors who use an agent to fill in this template:
%
%   Turn my notes below into a self-contained LaTeX problem using the attached
%   template (question-template.tex).
%   - Define every symbol; the statement must make sense without our Zulip.
%   - Keep labels honest: conjectures are conjectures, proposed directions are
%     not theorems. Do not invent results, citations or theorem numbers; write
%     [CHECK] wherever you are unsure.
%   - Give variants with stronger or weaker assumptions where natural.
%   - Fill in "What is known" only from sources I gave you, and credit
%     unpublished work by name.
%   - Make sure the file compiles with latexmk -pdf.
%   I will read and check the result before submitting it.
% ---------------------------------------------------------------------------
